\section{Warunki Hilberta--Bernaysa--Löba}\label{sec:warunki-hbl}
Do dowodu drugiego twierdzenia o niezupełności potrzebujemy własności predykatu dowodliwości, które sama reprezentowalność relacji dowodu nie zapewnia. W tym rozdziale $T$ jest efektywnie aksjomatyzowalnym rozszerzeniem PA, a $\operatorname{Bew}_T$ oznacza predykat z Lematu~\ref{res:7:7}. Piszemy $\Box\varphi$ zamiast $\operatorname{Bew}_T(\ulcorner\varphi\urcorner)$; w szczególności $\Box\Box\varphi$ oznacza $\operatorname{Bew}_T(\ulcorner\operatorname{Bew}_T(\ulcorner\varphi\urcorner)\urcorner)$.

\begin{definicja}[Warunki wyprowadzalności]\label{res:11:1}
Predykat $\Box$ spełnia \emph{warunki Hilberta--Bernaysa--Löba}, jeśli dla wszystkich zdań $\varphi,\psi$ zachodzą
\begin{align}
T\vdash\varphi&\quad\Longrightarrow\quad T\vdash\Box\varphi,\label{eq:11:1}\\
T&\vdash\Box(\varphi\to\psi)\to(\Box\varphi\to\Box\psi),\label{eq:11:2}\\
T&\vdash\Box\varphi\to\Box\Box\varphi.\label{eq:11:3}
\end{align}
Warunki te oznaczamy odpowiednio (D1)--(D3).
\end{definicja}

\begin{uwaga}\label{res:11:2}
(D1) jest regułą \emph{metateoretyczną}: z rzeczywistego dowodu $\varphi$ w $T$ tworzymy dowód $\Box\varphi$. Natomiast (D2) i (D3) są zdaniami \emph{dowodliwymi wewnątrz} $T$, także gdy kwantyfikowane kody dowodów mają niestandardowe wartości w modelach $T$. Dlatego (D3) nie wynika z samego (D1).
\end{uwaga}

\begin{lemat}[Operacje p.r. wewnątrz PA]\label{res:11:3}
Każda funkcja pierwotnie rekurencyjna $f:\N^k\to\N$ ma opis $F_f(\vec x,y)$ klasy $\Sigma_1$ taki, że
$\PA\vdash\forall\vec x\,\exists!y\,F_f(\vec x,y)$.
PA dowodzi ponadto równań definiujących funkcję $f$. Można więc stosować jej symbol jako skrót w rozszerzeniu definicyjnym PA; wszystkie takie skróty dają się wyeliminować na rzecz formuł języka arytmetyki.
\end{lemat}
\begin{proof}
W Lemacie~\ref{res:7:4} dowiedliśmy w metateorii, że każdy skończony ciąg ma kod $\beta$. Teraz potrzebujemy \emph{jednego dowodu w PA}, działającego także dla długości będącej zmienną. Najpierw, bez odwoływania się do niniejszego lematu, wykazujemy
\begin{equation}\label{eq:11:4}
\PA\vdash\forall a\,\forall b\,\forall n\,\forall c\,\exists a'\,\exists b'\,
\bigl[(\forall i<n\ \beta(a',b',i)=\beta(a,b,i))
\land\beta(a',b',n)=c\bigr].
\end{equation}
Tutaj $\beta$ jest skrótem definiowalnym przez formułę $B$ z Lematu~\ref{res:7:4}. Jest on dozwolony już na tym etapie, gdyż PA dowodzi istnienia i jednoznaczności reszty z dzielenia przez dodatni moduł, co ustalono niezależnie w dowodzie Lematu~\ref{res:7:5}.

\emph{Jednoczesne kodowanie zmiennej liczby reszt.} W PA potrafimy, przez indukcję po $n$, wskazać dodatnią wspólną wielokrotność $1,\ldots,n+1$ większą od zadanej liczby. Wybieramy więc $b'>a+c$ podzielne przez $1,\ldots,n+1$ i kładziemy $m_i=1+(i+1)b'$ dla $i\le n$. Rozumowanie z Lematu~\ref{res:7:4}, po zapisaniu w PA, dowodzi parami względnej pierwszości tych modułów. Określamy żądane reszty przez $\rho_i=\beta(a,b,i)$ dla $i<n$ oraz $\rho_n=c$; mamy $\rho_i<m_i$.

Aby zastosować chińskie twierdzenie o resztach \emph{wewnątrz PA}, nie wystarczy powołać się na jego prawdziwość w $\N$. PA dowodzi algorytmu Euklidesa i tożsamości Bézouta: przez indukcję uzasadnia dzielenie z resztą, a następnie formalizuje kolejne dzielenia; współczynniki całkowite koduje parami liczb naturalnych. Jeżeli $\gcd(u,v)=1$, otrzymujemy współczynniki $r,s$ spełniające $ru+sv=1$. Prowadzimy w PA indukcję po liczbie $k\le n+1$ uwzględnionych kongruencji. Jej niezmiennik stwierdza istnienie liczby $a_k$ spełniającej $a_k\equiv\rho_i\pmod{m_i}$ dla $i<k$ oraz wspólnego modułu $M_k=\prod_{i<k}m_i$. Przy dopisywaniu warunku $a\equiv\rho_k\pmod{m_k}$ korzystamy z $\gcd(M_k,m_k)=1$ i z tożsamości Bézouta, po czym zastępujemy uzyskaną liczbę jej resztą modulo $M_km_k$. Jest to indukcja po liczbie warunków, nie po \emph{już istniejącym} kodzie wynikowego ciągu. Dla $k=n+1$ otrzymujemy $a'$ spełniające wszystkie kongruencje, a tym samym~\eqref{eq:11:4}.

\emph{Kodowanie obliczeń zmiennej długości.} Z~\eqref{eq:11:4} i indukcji w PA wynika, że można kolejno dopisywać wartości do zakodowanego ciągu, sklejać dwa ciągi i wybierać zakodowaną rodzinę świadków. Na przykład dla ustalonej formuły $\theta(i,y)$ PA dowodzi
\[
\forall n\,\Bigl[(\forall i<n\ \exists y\,\theta(i,y))
\to\exists a\,\exists b\ \forall i<n\,\theta(i,\beta(a,b,i))\Bigr].
\]
Dowód jest indukcją po $n$: w kroku wybieramy świadka dla indeksu $n$ i używamy~\eqref{eq:11:4}. Tak otrzymujemy wewnątrz PA skończone, ale dowolnie długie ślady obliczeń, bez nieuprawnionego przejścia od osobnych dowodów dla numerałów do zdania ogólnego.

\emph{Indukcja po konstrukcji funkcji.} Dla funkcji bazowych PA dowodzi totalności i równań bezpośrednio. Przy składaniu wykorzystujemy opisy $F_g,F_{h_i}$ i dowiedzione w PA istnienie ich wartości; jednoznaczność funkcji składowych daje jednoznaczność wyniku. W kroku rekursji prostej
$f(\vec x,0)=g(\vec x)$,
$f(\vec x,S(n))=h(\vec x,n,f(\vec x,n))$
stosujemy indukcję PA po $n$. Hipoteza indukcyjna dostarcza zakodowany ciąg kolejnych wyników do chwili $n$; totalność $h$ pozwala wyznaczyć wynik następny, a~\eqref{eq:11:4} dołączyć go do kodu. Deterministyczność równań rekurencji daje jednoznaczność. Formułę $F_f$ wybieramy jako egzystencjalny, a zatem $\Sigma_1$, opis odpowiedniego skończonego śladu. Tę samą indukcję stosujemy do dowodów równań definiujących $f$.

Potrzebna poniżej jest jeszcze własność jednolita: dla ustalonego algorytmu realizującego funkcję p.r. możemy wybrać pierwotnie rekurencyjną procedurę tworzącą \emph{kod jego śladu obliczenia}. Dla funkcji bazowych ślady są jawne; dla składania są sklejane, a dla rekursji prostej budowane przez rekursję po jej argumencie. PA dowodzi poprawności tych konstrukcji tą samą indukcją. Nie otrzymujemy takiego dowodu totalności dla dowolnej funkcji obliczalnej z nieograniczoną minimalizacją.
\end{proof}

\begin{lemat}[Formalizacja łączenia dowodów]\label{res:11:4}
Dla ustalonych zdań $\varphi,\psi$ istnieje funkcja p.r. $m$ taka, że
\begin{equation}\label{eq:11:5}
\PA\vdash\forall p\,\forall q\,\bigl(
\operatorname{Prf}_T(p,\ulcorner\varphi\to\psi\urcorner)
\land\operatorname{Prf}_T(q,\ulcorner\varphi\urcorner)
\to\operatorname{Prf}_T(m(p,q),\ulcorner\psi\urcorner)\bigr).
\end{equation}
\end{lemat}
\begin{proof}
Dekodujemy pełne dowody $p,q$ z Lematu~\ref{res:7:7}. Za wierszami pierwszego umieszczamy wiersze drugiego, przesuwając ich odwołania do wcześniejszych wierszy o długość pierwszego dowodu. Dopisujemy $\psi$ z uzasadnieniem przez modus ponens, wskazując końcowe wiersze obu dowodów. Certyfikaty użytych aksjomatów przenosimy bez zmian, a do uzyskanej listy dołączamy ślad akceptacji przez ustalony weryfikator $V_T$. Wszystkie operacje na kodach oraz obliczenie śladu weryfikatora są p.r. na mocy Lematów~\ref{res:6:5} i~\ref{res:11:3}. Funkcję $m$ dopełniamy dowolną wartością na niepoprawnych wejściach.

Aby uzyskać~\eqref{eq:11:5}, musimy sprawdzić tę konstrukcję w PA, nie tylko dla rzeczywistych dowodów. Indukcja po liczbie kopiowanych wierszy zachowuje niezmiennik: skopiowany prefiks spełnia reguły weryfikatora, a każde odwołanie do przesłanki zostało przesunięte zgodnie z pozycją w nowym dowodzie. Wiersze aksjomatów zachowują certyfikaty, a dla wierszy otrzymanych przez reguły wnioskowania poprawność wynika z ich syntaktycznych testów. Nowy ostatni wiersz przechodzi test modus ponens na mocy przesłanek~\eqref{eq:11:5}. PA potwierdza także, że dołączony ślad weryfikatora kończy się akceptacją, wykorzystując formalizację procedury p.r. z Lematu~\ref{res:11:3}. To dowodzi~\eqref{eq:11:5} dla dowolnych $p,q$ spełniających jej poprzednik.

Nie wymagamy, aby PA rozstrzygała wszystkie aksjomaty $T$: w wyjściowych kodach znajdują się już ich świadectwa pochodzenia z efektywnego wyliczenia.
\end{proof}

\begin{lemat}[Sformalizowana $\Sigma_1$-zupełność]\label{res:11:5}
Niech $\sigma(\vec x)$ będzie ustaloną formułą $\Sigma_1$, a $\operatorname{inst}_\sigma(\vec n)=\ulcorner\sigma(\overline{\vec n})\urcorner$ funkcją podstawiania numerałów. Wtedy
\begin{equation}\label{eq:11:6}
\PA\vdash\forall\vec x\,\bigl(\sigma(\vec x)
\to\operatorname{Bew}_T(\operatorname{inst}_\sigma(\vec x))\bigr).
\end{equation}
Dla zdania $\sigma$ krotka $\vec x$ jest pusta, a $\operatorname{inst}_\sigma$ jest stałą równą $\ulcorner\sigma\urcorner$.
\end{lemat}
\begin{proof}
\emph{Jednolite dowody dla formuł ograniczonych.} Ustalmy formułę $\Delta_0$ $\delta(\vec x)$. Skonstruujemy funkcje p.r. $b_\delta^+,b_\delta^-$ wytwarzające kody dowodów prawdziwej instancji $\delta(\overline{\vec n})$ lub negacji fałszywej instancji. Co istotne, w PA udowodnimy \emph{jednolity} warunek ich poprawności, mianowicie
\begin{align}
\PA&\vdash\forall\vec x\,\bigl(\delta(\vec x)\to
\operatorname{Prf}_T(b_\delta^+(\vec x),\operatorname{inst}_\delta(\vec x))\bigr),\label{eq:11:7}\\
\PA&\vdash\forall\vec x\,\bigl(\neg\delta(\vec x)\to
\operatorname{Prf}_T(b_\delta^-(\vec x),\operatorname{inst}_{\neg\delta}(\vec x))\bigr).\label{eq:11:8}
\end{align}
Same dowody dla wszystkich ustalonych numerałów, bez twierdzeń~\eqref{eq:11:7}--\eqref{eq:11:8}, nie wystarczyłyby do wprowadzenia zmiennych świadków.

Konstrukcję prowadzimy przez indukcję w metateorii po budowie ustalonej formuły $\delta$; w każdym kroku uzyskujemy zarówno procedurę p.r., jak i dowód w PA jej poprawności dla wszystkich wartości parametrów.
\begin{enumerate}[label=(\alph*)]
\item \emph{Termy i atomy.} Rekurencyjnie obliczamy wartości termów przy podstawieniu numerałów i produkujemy ich formalne obliczenia z aksjomatów (Q4)--(Q7). Gdy dwie wartości są równe, łączymy dowody ich obliczenia; gdy są różne, stosujemy (Q1)--(Q2) do numerałów. Relację $<$ obsługujemy za pomocą (Q8) i skończonych odcinków z Lematu~\ref{res:4:4}. Operacje tworzenia kodów i dowodów są p.r.; PA dowodzi ich poprawności, korzystając z równań obliczeń i indukcji po liczbie zastosowań następnika.
\item \emph{Spójniki.} Z dowodów dla podformuł konstruujemy odpowiedni dowód dla $\neg,\land,\lor,\to$, korzystając z ustalonych tautologii i reguły modus ponens. Wybór właściwej gałęzi zależy od obliczonej wartości logicznej. Poprawność konstrukcji w PA wynika z poprawności funkcji dla podformuł oraz z formalizacji łączenia dowodów (Lemat~\ref{res:11:4}).
\item \emph{Kwantyfikatory ograniczone.} Dla formuły $\exists u<t(\vec x)\,\eta(\vec x,u)$ obliczamy $k=t(\overline{\vec n})$ i sprawdzamy kolejno $u=0,\ldots,k-1$. Jeśli formuła jest prawdziwa, wskazujemy świadka i dołączamy wprowadzenie kwantyfikatora; jeśli jest fałszywa, łączymy dowody negacji wszystkich instancji. Dla $\forall u<t\,\eta$ postępujemy dualnie. Przejście od skończonej koniunkcji lub alternatywy do formuły ograniczonej opieramy na dowodzie skończonego odcinka z Lematu~\ref{res:4:4}. Przy $k=0$ otrzymujemy dowód odpowiedniej pustej koniunkcji bądź zaprzeczenia pustej alternatywy.
\end{enumerate}

W ostatnim przypadku nie wolno ograniczyć się do instrukcji „połącz $k$ dowodów”. W PA stosujemy indukcję \emph{po zmiennej} $k$. Niezmiennik stwierdza, że skonstruowany kod zawiera poprawne uzasadnienia dotychczas sprawdzonych instancji, a jego końcowy wiersz odpowiada ich koniunkcji lub alternatywie, stosownie do rozpatrywanego przypadku. Indukcja używa jednolitego twierdzenia o poprawności procedury dla podformuły $\eta$, Lematu~\ref{res:11:4} oraz dowodliwej w PA poprawności operacji na kodach z Lematu~\ref{res:11:3}. Ponieważ głębokość składniowa $\delta$ jest ustalona, cała procedura pozostaje p.r. Dla funkcji $b_\delta^\pm$ wytwarzamy również końcowy certyfikat uruchomienia weryfikatora, aby kod spełniał dokładnie definicję $\operatorname{Prf}_T$ z Lematu~\ref{res:7:7}. Występujące w tych dowodach aksjomaty $Q$ zastępujemy ustalonymi $T$-dowodami ich instancji wraz z certyfikatami; istnienie takich skończonych dowodów wynika z $T\supseteq\PA\supseteq Q$. Tak otrzymujemy~\eqref{eq:11:7}--\eqref{eq:11:8} jako twierdzenia PA, a nie tylko rodzinę twierdzeń o poszczególnych numerałach.

\emph{Przejście do formuł $\Sigma_1$.} Zapiszmy $\sigma(\vec x)=\exists\vec u\,\delta(\vec x,\vec u)$, gdzie $\delta$ jest $\Delta_0$. Procedura p.r. bierze kod dowodu $\delta(\overline{\vec n},\overline{\vec k})$ otrzymany z $b_\delta^+$, a następnie dopisuje wprowadzenia kwantyfikatorów egzystencjalnych, otrzymując dowód $\sigma(\overline{\vec n})$. Na mocy Lematu~\ref{res:11:3} PA dowodzi poprawności tej operacji, zatem dla odpowiedniej funkcji p.r. $c$ mamy
$\PA\vdash\forall\vec x\,\forall\vec u\,\bigl(\delta(\vec x,\vec u)\to\operatorname{Prf}_T(c(\vec x,\vec u),\operatorname{inst}_\sigma(\vec x))\bigr)$.
Wprowadzając kwantyfikator po kodzie dowodu, a następnie po świadkach $\vec u$, otrzymujemy~\eqref{eq:11:6}.
\end{proof}

\begin{lemat}[Warunki wyprowadzalności dla standardowego predykatu]\label{res:11:6}
Dla każdej efektywnie aksjomatyzowalnej teorii $T\supseteq\PA$ predykat $\operatorname{Bew}_T$ z Lematu~\ref{res:7:7} spełnia warunki (D1)--(D3).
\end{lemat}
\begin{proof}
\emph{(D1).} Jeżeli $T\vdash\varphi$, to istnieje konkretny pełny kod $p$ dowodu $\varphi$. Z Lematu~\ref{res:7:7} mamy $Q\vdash\operatorname{Prf}_T(\overline p,\ulcorner\varphi\urcorner)$; wprowadzenie kwantyfikatora egzystencjalnego daje $T\vdash\Box\varphi$.

\emph{(D2).} Stosujemy~\eqref{eq:11:5} i wprowadzamy w PA kwantyfikatory egzystencjalne po kodach $p,q$. Otrzymujemy
$\PA\vdash\Box(\varphi\to\psi)\land\Box\varphi\to\Box\psi$, a więc także implikację (D2) w $T$.

\emph{(D3).} Formuła $\sigma:=\operatorname{Bew}_T(\ulcorner\varphi\urcorner)$ jest zdaniem $\Sigma_1$. Z Lematu~\ref{res:11:5} wynika
$\PA\vdash\operatorname{Bew}_T(\ulcorner\varphi\urcorner)\to
\operatorname{Bew}_T(\ulcorner\operatorname{Bew}_T(\ulcorner\varphi\urcorner)\urcorner)$,
czyli (D3).
\end{proof}

\begin{uwaga}[Dlaczego wybór predykatu ma znaczenie]\label{res:11:7}
Dwie formuły mogą definiować tę samą relację na liczbach standardowych, choć PA nie dowodzi ich równoważności. Warunki (D2)--(D3) odnoszą się do konkretnej reprezentacji predykatu dowodliwości oraz dowodliwej wewnątrz PA poprawności operacji na kodach. Nie wolno przypisywać ich dowolnemu predykatowi reprezentującemu zbiór twierdzeń ani automatycznie przenosić na predykat Rossera.
\end{uwaga}

